\documentclass[acmtog,anonymous,review]{acmart}
\acmSubmissionID{1107}

\usepackage{bm}
\usepackage{mathtools}
\usepackage{enumitem}

\newcommand{\norm}[1]{\left\lVert #1 \right\rVert}
\newcommand{\abs}[1]{\left\lvert #1 \right\rvert}
\newcommand{\R}{\mathbb{R}}
\newcommand{\T}{\mathsf{T}}

\DeclareMathOperator{\diag}{diag}

\title{Supplementary Material:\\
Continuous Collision Detection Trust Regions for Clothing Simulation}

\begin{document}
\maketitle

\subsection{Incremental Potential and Energy Decomposition}
\label{sec:incremtal_energy}

We adopt the standard incremental potential
formulation~\cite{li2020ipc}. A backward Euler step reduces to the
minimization
\begin{equation}
\label{eq:inc_pot}
\Phi(\mathbf{x})
=
\tfrac{1}{2}\,
\|\mathbf{x}-\hat{\mathbf{x}}\|_{\mathbf{M}}^{2}
+
(\Delta t)^2\,PE(\mathbf{x}),
\end{equation}
where
$\|\mathbf{v}\|_{\mathbf{M}}^{2}
=\mathbf{v}^{\mathsf T}\mathbf{M}\mathbf{v}$,
$\hat{\mathbf{x}}=\mathbf{x}^n+\Delta t\,\mathbf{v}^n$ is the
inertial extrapolation, and
$\mathbf{M}=\diag(m_1\mathbf{I}_3,\dots,m_N\mathbf{I}_3)$
is the lumped mass matrix. We decompose the total potential as
\begin{multline}
\label{eq:pe_decomp}
PE(\mathbf{x})
=
PE^{\mathrm{e}}(\mathbf{x})
+
PE^{\mathrm{b}}(\mathbf{x})
+
PE^{\mathrm{g}}(\mathbf{x})
\\
+\;
PE^{\mathrm{bar}}(\mathbf{x})
+
PE^{\mathrm{pin}}(\mathbf{x}),
\end{multline}
corresponding respectively to corotated elastic stretching,
discrete shell bending, gravity, IPC barrier contact, and pinning
constraints. We now briefly state each term.

\paragraph{Corotated elasticity.}
On each triangle with rest-shape matrix
$\mathbf{D}_m\in\mathbb{R}^{2\times 2}$, deformed segment matrix
$\mathbf{D}_s\in\mathbb{R}^{3\times 2}$, and deformation gradient
$\mathbf{F}=\mathbf{D}_s\mathbf{D}_m^{-1}$, we use the corotated
energy density~\cite{stomakhin2012elasticity}
\begin{equation}
\label{eq:psi_corotated}
\psi(\mathbf{F})
=
\mu\|\mathbf{F}-\mathbf{R}\|_F^2
+
\tfrac{\lambda}{2}(J-1)^2,
\end{equation}
where $\mathbf{F}=\mathbf{R}\mathbf{S}$ is the polar decomposition,
$J=\sqrt{\det(\mathbf{F}^{\mathsf T}\mathbf{F})}$, and $\mu,\lambda$
are Lam\'{e} parameters. The element elastic energy is
$E_t=A_t\,\psi(\mathbf{F})$ with rest area
$A_t=\tfrac{1}{2}|\det\mathbf{D}_m|$, and the total stretching
energy is $PE^{\mathrm{e}}=\sum_t E_t$.

\paragraph{Bending.}
Following Grinspun et al.~\cite{grinspun2003discrete}, the discrete
shell bending energy for a hinge segment $e$ with dihedral-angle
complement $\theta$ and rest value $\bar{\theta}$ is
\begin{equation}
\label{eq:bending_energy}
w_e = k_B\,c_e\,(\theta - \bar{\theta})^2,
\qquad
c_e = \frac{|\bar{e}|}{\bar{h}_e},
\end{equation}
where $k_B$ is the bending stiffness and $c_e$ is a geometric
coefficient depending on the rest segment length and opposite heights.
The total bending energy is $PE^{\mathrm{b}}=\sum_e w_e$.

\paragraph{Gravity and pinning.}
Gravity contributes a linear potential
$PE_i^{\mathrm{g}}(\mathbf{x}_i)
=-m_i\,\mathbf{g}_{\mathrm{grav}}^{\mathsf T}\mathbf{x}_i$
at each node. Pinned nodes are enforced through a quadratic penalty
$PE_i^{\mathrm{pin}}(\mathbf{x}_i)
=\tfrac{k_{\mathrm{pin}}}{2}
\|\mathbf{x}_i-\mathbf{x}_i^{\mathrm{pin}}\|^2$
with $k_{\mathrm{pin}}\gg 1$.

\paragraph{IPC barrier.}
Non-penetration between deformable primitives is enforced by the IPC
log-barrier~\cite{li2020ipc}. For a primitive pair with gap
$\delta>0$, the barrier function is
\begin{equation}
\label{eq:barrier}
b(\delta;\hat{d})
=
\begin{cases}
-(\delta-\hat{d})^2\ln(\delta/\hat{d}), & 0<\delta<\hat{d},\\
0, & \delta\ge\hat{d},
\end{cases}
\end{equation}
and the total barrier energy is
\begin{equation}
\label{eq:barrier_energy}
PE^{\mathrm{bar}}(\mathbf{x}) = \kappa_{\mathrm{bar}} \sum_{(p,q)} b(\delta_{pq};\hat{d}),
\end{equation}
where $\kappa_{\mathrm{bar}}$ is the barrier stiffness and the sum
ranges over all active vertex--triangle and edge--edge pairs.

\subsection{Signed-Distance Penalty Energy}

We additionally enforce non-penetration with respect to rigid
obstacles through a penalty energy composed from the signed distance
function and a piecewise-linear Heaviside. Let
$\Omega \subset \mathbb{R}^3$ denote the obstacle region with signed
distance function $\phi(\mathbf{x})$, negative inside and positive
outside. Define the piecewise-linear Heaviside
\[
H(z) =
\begin{cases}
1, & z < 0, \\
(\epsilon - z) / \epsilon, & 0 \le z \le \epsilon, \\
0, & z > \epsilon,
\end{cases}
\]
so that $H$ transitions linearly from full penetration to zero over
a band of width $\epsilon$. The per-node penalty energy is
\begin{equation}
\label{eq:pe_sdf}
PE_i^{\mathrm{sdf}}(\mathbf{x}_i)
= \tfrac{k_{\mathrm{sdf}}}{2}
\bigl[H\bigl(\phi(\mathbf{x}_i)\bigr)\bigr]^2,
\end{equation}
where $k_{\mathrm{sdf}}$ is a stiffness parameter. Nodes with
$\phi(\mathbf{x}_i) > \epsilon$ contribute zero energy and are
skipped. 
\subsection{Local Newton Subproblem}
\label{sec:local_newton}

The Gauss--Seidel strategy updates nodes sequentially through local
implicit subproblems. Each local step incorporates the inertia,
elastic, gravitational, pinning, and active barrier contributions
associated with the current node.

For node $i$, the local first-order optimality condition is
\begin{multline}
\label{eq:local_opt}
\nabla_{\mathbf{x}_i}\Phi(\mathbf{x})
=
m_i(\mathbf{x}_i-\hat{\mathbf{x}}_i)
+
(\Delta t)^2
\bigl(
\nabla_{\mathbf{x}_i}PE^{\mathrm{e}}
\\
+\nabla_{\mathbf{x}_i}PE^{\mathrm{b}}
+\nabla_{\mathbf{x}_i}PE^{\mathrm{bar}}
+\nabla_{\mathbf{x}_i}PE^{\mathrm{pin}}
-m_i\mathbf{g}_{\mathrm{grav}}
\bigr)
=\mathbf{0},
\end{multline}
where the pinning term is present only if node $i$ is pinned.
Linearizing about the current $\mathbf{x}_i$ and solving for the
correction $\boldsymbol{\delta}_i\in\mathbb{R}^3$ gives
\begin{equation}
\label{eq:local_newton}
\mathbf{H}_{ii}\,\boldsymbol{\delta}_i = \mathbf{g}_i,
\end{equation}
where $\mathbf{g}_i=\nabla_{\mathbf{x}_i}\Phi(\mathbf{x})$ and
\begin{multline}
\label{eq:local_hessian}
\mathbf{H}_{ii}
=
m_i\mathbf{I}_3
+
(\Delta t)^2
\bigl(
\mathbf{H}^{\mathrm{e}}_{ii}
+\mathbf{H}^{\mathrm{b}}_{ii}
\\
+\mathbf{H}^{\mathrm{bar}}_{ii}
+\mathbf{H}^{\mathrm{pin}}_{ii}
\bigr)
\in\mathbb{R}^{3\times 3}.
\end{multline}
Here $\mathbf{H}^{\mathrm{e}}_{ii}$ is the diagonal block of the
elastic Hessian summed over all triangles incident to node $i$,
$\mathbf{H}^{\mathrm{b}}_{ii}$ is the diagonal block of the bending
Hessian summed over all hinge segments incident to node $i$,
$\mathbf{H}^{\mathrm{bar}}_{ii}$ is the diagonal block of the barrier
Hessian summed over all active barrier pairs involving node $i$, and
$\mathbf{H}^{\mathrm{pin}}_{ii}=k_{\mathrm{pin}}\mathbf{I}_3$ if $i$
is pinned and zero otherwise. The node position is then updated as
\begin{equation}
\label{eq:node_update}
\mathbf{x}_i
\leftarrow
\mathbf{x}_i - \omega_i\,\boldsymbol{\delta}_i,
\end{equation}
where $\omega_i\in(0,1]$ is a step size determined by
linear CCD in
Section~\ref{sec:linear_ccd}.

\subsection{Initial Guess}
\label{sec:initial_guess}

The initial iterate $\mathbf{x}^{(0)}$ must be collision-free and
connected to $\mathbf{x}^n$ by a linear trajectory so that CCD
queries are well defined.

The inertial term
$\tfrac{1}{2}\|\mathbf{x}-\hat{\mathbf{x}}\|_{\mathbf{M}}^2$ in
the incremental potential~\eqref{eq:inc_pot} is minimized at the inertial
extrapolation
$\hat{\mathbf{x}}=\mathbf{x}^n+\Delta t\,\mathbf{v}^n$, which is
typically the dominant contribution to $\Phi$. A natural initial
guess is therefore to advance toward $\hat{\mathbf{x}}$ as far as
collision constraints allow. Since all nodes move simultaneously
along this trajectory, the coplanarity polynomial for each candidate
primitive pair is a full cubic in
$t$~\cite{bridson2002cloth}. For every candidate pair we compute the
earliest collision time $t^*\in[0,1]$ and set
\begin{equation}
\label{eq:omega0}
\omega_0
=
\min \,\!\bigl(1,\;
\min_{\text{all pairs}}\,0.9\,t^*\bigr),
\end{equation}
\begin{equation}
\label{eq:ccd_guess}
\mathbf{x}^{(0)}
=
\mathbf{x}^n+\omega_0\,\Delta t\,\mathbf{v}^n.
\end{equation}
This places the initial guess at the collision-free configuration
that minimizes the inertial energy along the line from
$\mathbf{x}^n$ to $\hat{\mathbf{x}}$, so the solver only needs to
resolve the potential energies.

\subsection{Linear CCD}
\label{sec:linear_ccd}

Each local Newton update must be filtered to ensure collision freedom.
We use CCD based on coplanarity
tests~\cite{bridson2002cloth} to determine the maximal admissible
step size.  Candidate primitive pairs are identified via swept
axis-aligned bounding boxes organized in bounding volume
hierarchies~\cite[Section~4.2]{ericson2004collision}.

Since only a single node moves per local step, the coplanarity
polynomial simplifies significantly.  Given four points interpolating
as $\mathbf{x}_i(t)=\mathbf{x}_i^0+t\,\mathbf{d}_i$ with
$t\in[0,1]$, define the edge vectors
$\mathbf{p}=\mathbf{x}_1-\mathbf{x}_0$,
$\mathbf{q}=\mathbf{x}_2-\mathbf{x}_0$,
$\mathbf{r}=\mathbf{x}_3-\mathbf{x}_0$.  The scalar triple product
$f(t)=(\mathbf{p}(t)\times\mathbf{q}(t))\cdot\mathbf{r}(t)$ is
cubic in general, but when only one vertex carries nonzero
displacement, at most one of $\mathbf{p},\mathbf{q},\mathbf{r}$
depends on~$t$, so
\begin{equation}\label{eq:linear_ccd}
  f(t) = ct + d, \qquad t^* = -d/c.
\end{equation}
For example, when the query vertex~$\mathbf{x}_0$ moves with
displacement~$\mathbf{d}_0$, the coefficients are
$d=(\mathbf{p}_0\!\times\!\mathbf{q}_0)\cdot\mathbf{r}_0$ and
$c=-(\mathbf{p}_0\!\times\!\mathbf{q}_0)\cdot\mathbf{d}_0
   +(\mathbf{p}_0\!\times\!\mathbf{d}_0)\cdot\mathbf{r}_0
   +(\mathbf{d}_0\!\times\!\mathbf{q}_0)\cdot\mathbf{r}_0$,
where $\mathbf{p}_0,\mathbf{q}_0,\mathbf{r}_0$ are the edge vectors
at~$t\!=\!0$.  A candidate~$t^*$ is accepted only if $|c|>\epsilon$,
$t^*\!\in\!(0,1]$, and the contact point lies inside the primitive
(positive barycentric coordinates for vertex--triangle; both edge
parameters in~$[0,1]$ for edge--edge).  The step is clamped to
$\hat{t}=0.9\,\min t^*$ to maintain a safety margin.

\bibliographystyle{ACM-Reference-Format}
\bibliography{refs}
\end{document}